\input{../tex-inputs/latex-leseansicht-vorspann}

\section[Arthur Schnitzler an Theodor von Sosnosky, 10. 10. 1901]{L04235 Arthur Schnitzler an Theodor von Sosnosky, 10. 10. 1901}
\nopagebreak\mylabel{L04235v}
\rehead{ }\normalsize\beginnumbering\briefempfaengerindex{Sosnosky, Theodor von@\textsc{Sosnosky, Theodor von}!zzzSchnitzler, Arthur@\emph{von Arthur Schnitzler}!1901-10-101@{10. 10. 1901}|(be}
\toendnotes[C]{\smallbreak\pagebreak[2]}
\correspDesc{Versand  durch Arthur Schnitzler am 10. 10. 1901 in Wien
\newline{}Übermittlung  am 11. 10. 1901 in Wien 9
\newline{}Erhalt  durch Theodor von Sosnosky am 12. 10. 1901 in Alland}\toendnotes[C]{\smallbreak}
\Standort{CUL, Schnitzler, B 126.}
\physDesc{Brief, 2 Blätter, 7 Seiten, Kuvert, 2632 Zeichen
\newline{}Handschrift: schwarze Tinte, deutsche Kurrent
\newline{}Versand: 1) Stempel: »\nobreak{}\oindex{XXXX Ortsangabe fehlt|pwk}Wien 9/3 72, 11. 10. 0\textcolor{gray}{1}, 10–11\textcolor{gray}{V}\nobreak{}«.   2) Stempel: »\nobreak{}\oindex{XXXX Ortsangabe fehlt|pwk}Alland, 12 10 {[}01{]}\nobreak{}«. 
\newline{}Zusatz: Die beiden Briefe Schnitzlers an Sosnosky wurden (mutmaßlich von Heinrich Schnitzler\pwindex{\textcolor{red}{\textsuperscript{XXXX indx1}}|pw}) auf
                                 einer Versteigerung im Wiener Dorotheum am 27. 5. 1931 um
                                 35 Schilling erworben. Bleistiftvermerke: »5469«, »20–« und »K63a« }
\buchAbdrucke{\weitereDrucke{1) Theodor v. Sosnosky: \emph{Unveröffentlichte Schnitzler-Briefe über die »Leutnant-Gustl«-Affäre. Eine Sensation vor dreißig Jahren.} In: \emph{Neues Wiener Journal}, Jg. 39, Nr. 13.624, 26. 10. 1931, S. 4.} \weitereDrucke{2) Arthur Schnitzler: \emph{Lieutenant Gustl}. Herausgegeben und kommentiert von Ursula Renner unter Mitarbeit von Heinrich Bosse. Frankfurt am Main: \emph{Suhrkamp} 2007, S. 54 (Suhrkamp Basis
                        Bibliothek, 33).} \weitereDrucke{3) Ursula Renner: \emph{Dokumentation eines Skandals Arthur Schnitzlers »Lieutenant Gustl«.} In: \emph{Hofmannsthal-Jahrbuch}, Jg. 15 (2007), S. 210–211.} }\toendnotes[C]{\smallbreak}\pstart{}{\pb}Herrn Theodor von Sosnosky\pend{}\pstart{}\textsc{Alland bei Baden}\oindex{XXXX Ortsangabe fehlt|pw}\pend{}\pstart{}(bei Wien\oindex{Wien@\textbf{Wien}, \emph{Verwaltungsgebiet}|pw})\pend{}{\bigskip}\vspace{1em}
\pstart\center{}{\pb}Sehr geehrter Herr von Sosnosky,\pend\vspace{0.5em}
\pstart
           ich danke Ihnen verbindlichſt für die Überſendung Ihres Artikels\pwindex{Sosnosky, Theodor von 4.\,1.\,1866 Budapest – 13.\,2.\,1943 Wien@\textsc{Sosnosky, Theodor von} (4.\,1.\,1866 Budapest – 13.\,2.\,1943 Wien), \emph{Schriftsteller}!Fall Schnitzler. Eine unbefangene Betrachtung@\strich\emph{Der Fall Schnitzler. Eine unbefangene Betrachtung}|pwv} über den Lieutenant Guſtl\pwindex{Schnitzler, Arthur 15.\,5.\,1862 Wien – 21.\,10.\,1931 ebd.@\textsc{Schnitzler, Arthur} (15.\,5.\,1862 Wien – 21.\,10.\,1931 ebd.), \emph{Schriftsteller, Mediziner}!Lieutenant Gustl. Novelle@\strich\emph{Lieutenant Gustl. Novelle}|pw}. Das literarische Leben wäre wahrlich eine
               angenehme und reinliche Sache, wenn man nur mit{ }ſo vornehmen Gegnern zu thun hätte
               wie Sie einer{ }ſind. Über die Novelle\pwindex{Schnitzler, Arthur 15.\,5.\,1862 Wien – 21.\,10.\,1931 ebd.@\textsc{Schnitzler, Arthur} (15.\,5.\,1862 Wien – 21.\,10.\,1931 ebd.), \emph{Schriftsteller, Mediziner}!Lieutenant Gustl. Novelle@\strich\emph{Lieutenant Gustl. Novelle}|pwv}{ }ſelbst wollen wir nicht mehr reden; hier iſt eine Verſtändigung
               unmöglich – denn Sie{ }ſind {\pb}Sie und ich bin ich. Hingegen muſs ich, was Sie vielleicht
               überraſchen wird, den Ehrenrath\orgindex{[Ehrenrat der österreichisch-ungarischen Armee zur Veröffentlichung von Lieutenant Gustl, 1901]@[Ehrenrat der österreichisch-ungarischen Armee zur Veröffentlichung von Lieutenant Gustl, 1901]|pw}, der mich{ }ſchuldig geſprochen, in Schutz nehmen. In dem Urtheil heißt es klar und deutlich,
               daſs ich die Standesehre \strikeout{dadurch} verletzt, \strikeout{da} indem ich durch meine Novelle\pwindex{Schnitzler, Arthur 15.\,5.\,1862 Wien – 21.\,10.\,1931 ebd.@\textsc{Schnitzler, Arthur} (15.\,5.\,1862 Wien – 21.\,10.\,1931 ebd.), \emph{Schriftsteller, Mediziner}!Lieutenant Gustl. Novelle@\strich\emph{Lieutenant Gustl. Novelle}|pwv} das Anſehn und die Ehre der oeſterreichiſch-ungariſchen Armee\orgindex{Gemeinsame Armee der österreichisch-ungarischen Doppelmonarchie@Gemeinsame Armee der österreichisch-ungarischen Doppelmonarchie|pw} geſchädigt und
               herabgeſetzt habe. Hätten die Herren gefunden, daſs ich die Standesehre {\pb}dadurch
               verletzt habe, daſs ich nicht perſönlich vor \strikeout{I} ihnen
               erſchienen bin, so hätten{ }ſie das{ }ſelbstverständlich erwähnt, und es geht kaum an,
               ihnen zuzumuthen, daſs{ }ſie durch mein perſönliches Erſcheinen beſtimmt worden wären,
               ein andres Urtheil zu fällen als{ }ſie nach dem Studium meiner Novelle\pwindex{Schnitzler, Arthur 15.\,5.\,1862 Wien – 21.\,10.\,1931 ebd.@\textsc{Schnitzler, Arthur} (15.\,5.\,1862 Wien – 21.\,10.\,1931 ebd.), \emph{Schriftsteller, Mediziner}!Lieutenant Gustl. Novelle@\strich\emph{Lieutenant Gustl. Novelle}|pwv} zu fällen für richtig hielten, da
               ja dieſe Novelle\pwindex{Schnitzler, Arthur 15.\,5.\,1862 Wien – 21.\,10.\,1931 ebd.@\textsc{Schnitzler, Arthur} (15.\,5.\,1862 Wien – 21.\,10.\,1931 ebd.), \emph{Schriftsteller, Mediziner}!Lieutenant Gustl. Novelle@\strich\emph{Lieutenant Gustl. Novelle}|pwv} nun einmal
               geſchrieben und veröffent{\pb}licht \strikeout{war} – und damit die
               von der Anklage angeno{\geminationm}ene Verletzung der Standesehre begangen war. Daſs ich mich
               \strikeout{wegen} entschuldigen würde, meine Novelle\pwindex{Schnitzler, Arthur 15.\,5.\,1862 Wien – 21.\,10.\,1931 ebd.@\textsc{Schnitzler, Arthur} (15.\,5.\,1862 Wien – 21.\,10.\,1931 ebd.), \emph{Schriftsteller, Mediziner}!Lieutenant Gustl. Novelle@\strich\emph{Lieutenant Gustl. Novelle}|pwv} geſchrieben zu haben, dürfte wohl
               von keiner Seite voraus geſetzt werden, und hätte ich mich dazu verſtanden,{ }ſo wäre
               meinem Empfinden nach darin – die einzige wirkliche Verletzung der\introOben{}jenigen\introOben{} Standesehre
               gelegen, die ich überhaupt anerkenne: das zu sagen \introOben{}und zu thun\introOben{}, was man {\pb}für richtig
               hält. – Was die zweite Anſchuldigung anbelangt, ich hätte gegen die perſönlichen
               Angriffe der Reichswehr\orgindex{Reichswehr@Reichswehr|pw} keine Schritte
               unternommen,{ }ſo theile ich Ihnen mit, daſs ich von dieſem Theil der Anklage erſt –
               aus dem Wortlaut des Urtheils Kenntnis erhielt. Was darüber zu{ }ſagen ist – mögen Sie
               in Ihrem eignen Artikel\pwindex{Sosnosky, Theodor von 4.\,1.\,1866 Budapest – 13.\,2.\,1943 Wien@\textsc{Sosnosky, Theodor von} (4.\,1.\,1866 Budapest – 13.\,2.\,1943 Wien), \emph{Schriftsteller}!Fall Schnitzler. Eine unbefangene Betrachtung@\strich\emph{Der Fall Schnitzler. Eine unbefangene Betrachtung}|pwv}
               nachleſen. Ich brauche Sie wohl nicht zu verſichern, daſs die Kränkung, die mir durch
               die \label{K_L04235-1v}\edtext{Enuncia{\pb}tionen}{\lemma{\textnormal{\emph{Enunciationen}}}\Cendnote{\textnormal{Enunziation: Aussage}}}\label{K_L04235-1} der feindlichen
               Preſſe verurſacht worden{ }ſind, etwa eben so tief gehen – wie die Freude,
               die mir ein Theil der zuſtimmenden Preßkundgebungen bereitet hat. Ich
               weiſs in Hinſicht auf die Einen wie die andern, was ich von ihrer Sachlichkeit wie
               Unpartheilichkeit zu halten habe. – Ihre Befürchtung, daſs es jetzt am Ende Leute
               geben könnte, die mir meinen Gruſs nicht erwidern, {\pb}iſt unbegründet; ich weiſs ganz
               gut, welchen Leuten ich die Ehre erweiſen darf,{ }ſie zuerst zu grüßen. –\pend
           
\pstart
           Nehmen Sie, verehrteſter Herr von Sosnosky, nochmals meinen Dank und die
               Verſicherung meiner aufrichten Hochachtung entgegen {\\[\baselineskip]}Ihr ergebener{\\[\baselineskip]}\spacefill\mbox{Arthur Schnitzler}\pend
           \leftskip=0em{}
\pstart
           Wien\oindex{Wien@\textbf{Wien}, \emph{Verwaltungsgebiet}|pw}, 10. 10. 901.\pend
           \selectlanguage{ngerman}\endnumbering\briefempfaengerindex{Sosnosky, Theodor von@\textsc{Sosnosky, Theodor von}!zzzSchnitzler, Arthur@\emph{von Arthur Schnitzler}!1901-10-101@{10. 10. 1901}|)be}\mylabel{L04235h}
\begin{anhang}
\end{anhang}\newcommand{\dateiname}{L04235}\newcommand{\titel}{Arthur Schnitzler an Theodor von Sosnosky, 10. 10. 1901}\newcommand{\editorInnen}{Herausgegeben von Jahnke, SelmaMüller, Martin Anton}\input{../tex-inputs/latex-leseansicht-abspann}
